\documentclass[../main]{subfiles}
\begin{document}
\chapter{Interacci\'on electromagn\'etica}
\info{Contenido}{
  Interacci\'on entre la corriente el\'ectrica y campo magn\'etico, generaci\'on de f.e.m y fuerzas.
  f.e.m. inducida por un conductor que se mueve en un campo magn\'etico.
}
\info{F\'ormulas}{
  \begin{center}
    \begin{tabular}{cc}
      $\epsilon = B l v = N \frac{\Delta \Phi}{\Delta t}$ & $\epsilon_{auto} = \frac{\Delta \Phi}{\Delta t}$ \\
      $e_{auto} = L \frac{\Delta I}{\Delta t}$ & $L = N \frac{\Phi}{I}$ \\
      $F = B l I$ & $\Phi = B s$
    \end{tabular}
  \end{center}

  Siendo:
  \begin{itemize}
    \item $s$: secci\'on transversal \unit{m^2}
    \item $\Phi$: Flujo magn\'etico \unit{Wb} (Weber = Tesla $\cdot m^2$)
    \item $\Delta \Phi = \Phi_{final} - \Phi_{inicial}$: Variaci\'on de flujo magn\'etico.
    \item $\Delta t = t_{final} - t_{inicial}$: Variaci\'on de tiempo \unit{s}
    \item $I$: Corriente el\'ectrica \unit{A}
    \item $\Delta I = I_{final} - I_{inicial}$: Variaci\'on de corriente \unit{A}
    \item $L$: Coeficiente de autoinducci\'on \unit{H} (Henrio)
    \item $N$: N\'umero de espiras.
    \item $F$: Fuerza electromotriz \unit{N}
    \item $\epsilon$: f.e.m. inducida \unit{V}
    \item $\epsilon_{auto}$: f.e.m. de autoinducci\'on \unit{V}
    \item $B$: Inducci\'on magn\'etica \unit{T}
    \item $l$: Longitud del conductor \unit{m}
    \item $v$: Velocidad del conductor \unit{m/s}
  \end{itemize}
}

\begin{tcolorbox}
  Teor\'ia
  \tcblower
  \begin{tabular}{|p{12cm}|c|}\hline
    Descripci\'on & C\'odigo \\ \hline
    f.e.m inducida por un conductor que se mueve en un campo magn\'etico & \qrcode{https://www.youtube.com/watch?v=6N7ukJxvQ_M}\\\hline
    Coeficiente de autoinducci\'on &
    \qrcode{https://www.youtube.com/watch?v=Xb5GHlaWv7g} \\\hline
    Fuerza producida en un conductor por el que circula corriente dentro de un campo magn\'etico &
    \qrcode{https://www.youtube.com/watch?v=QvBbUOn_T4k}\\ \hline
  \end{tabular}
\end{tcolorbox}

\actividad{Resolver los siguientes problemas}

\begin{enumerate}
  \item Un conductor se desplaza a una velocidad lineal de $v=\qty{10}{m/s}$ dentro de un campo magn\'etico fijo de $B=\qty{1,5}{T}$ de inducci\'on. Determinar el valor de la f.e.m. inducida en el mismo si posee una longitud de $l=\qty{1}{m}$.

  \item ¿A qu\'e velocidad se deber\'a mover un conductor de longitud $l=\qty{0,5}{m}$ para que dentro de una inducci\'on de campo magn\'etico de $B=\qty{1}{T}$ produzca una f.e.m de $e=\qty{5}{V}$?

  \item Determinar la inducci\'on de campo magn\'etico $B$, que ser\'a necesaria para que un conductor de $l=\qty{1}{m}$ de longitud que se mueve a $v=\qty{10}{m/s}$ cortando las l\'ineas del campo magn\'etico produzca una f.e.m de $e=\qty{12}{V}$.

  \item ¿Cu\'al la f.e.m de auto-inducci\'on si el flujo magn\'etico $\Phi$ que existe en el campo magn\'etico producido por una bobina de una electrov\'alvula si tiene esta un n\'ucleo cil\'indrico de di\'ametro $d=\qty{4}{cm}$ y la inducci\'on magn\'etica en la misma crece de $B=\qty{0}{T}$ a $B=\qty{1,5}{T}$ en $\Delta t=\qty{0,3}{s}$?

  \item Calcular el valor de la f.e.m de auto-inducci\'on que se desarrollar\'a en una bobina con coeficiente de auto-inducci\'on de $\qty{50}{mH}$ si se aplica una corriente que crece desde cero hasta $I=\qty{20}{A}$ en un tiempo de $\Delta t=\qty{0,1}{s}$.

  \item ¿Cu\'anto tiempo habr\'a que esperar para que la f.e.m de auto-inducci\'on se reduzca a $e_{auto}=\qty{1}{V}$ si se hace circular $I=\qty{10}{A}$ por una bobina con un coeficiente de auto-inducci\'on de $L=\qty{60}{mH}$?

  \item ¿Cu\'al es la f.e.m de auto inducci\'on del problema anterior cuando pasa \qty{1}{ms} de tiempo $t=\qty{0,001}{s}$ y ¿Cu\'al es el valor de la fem de auto-inducci\'on despu\'es de \qty{1}{min}?

  \item Con los datos del problema anterior elaborar un gr\'afico donde en el eje vertical (ordenadas) se coloque el valor de f.e.m. auto-inducci\'on $e_{auto}$ en el eje de las abscisas los valores de tiempo para $t_1=\qty{0,01}{s}$, $t_2=\qty{0,1}{s}$, $t_3=\qty{0,5}{s}$, $t_4=\qty{0,8}{s}$, $t_5=\qty{1}{s}$, $t_6=\qty{2}{s}$, $t_7=\qty{4}{s}, t_8=\qty{5}{s}$ y $t_9=\qty{10}{s}$.

  \item Una bobina de $N=200$ espiras se mueve cortando perpendicularmente un campo magn\'etico. La variaci\'on de flujo experimentado es uniforme y va desde los $\Phi_1=\qty{5}{mWb}$ a los $\Phi_2=\qty{10}{mWb}$ en un intervalo de tiempo de $\Delta t = \qty{0,5}{s}$. Averiguar la f.e.m inducida.

  \item Una bobina de $N=100$ espiras corta con un \'angulo de $45^\circ$ un campo magn\'etico uniforme, la variaci\'on de flujo experimentada es uniforme y es de $\Delta \Phi = \qty{4}{mWb}$ en $\Delta t = \qty{1}{s}$. Averiguar la f.e.m inducida.

  \item Los terminales de una bobina son conectados a una resistencia $R=\qty{1}{\Omega}$ y se genera una corriente de $I=\qty{0,5}{A}$ y cuando esta se mueve produce una variaci\'on de flujo de $\Delta \Phi =\qty{5}{mWb}$ en $\Delta t = \qty{2}{s}$ ¿Cu\'antas espiras N tiene?

  \item Hallar la longitud que tiene que tener un conductor que corta perpendicularmente un campo magn\'etico el cual posee una inducci\'on de $B=\qty{3}{T}$, y movi\'endose a una velocidad de $v=\qty{5}{m/s}$ genera una f.e.m de $e=\qty{10}{V}$.

  \item Calcular el valor de la f.e.m. de autoinducci\'on que desarrollar\'a una bobina con un coeficiente de autoinducci\'on de $L=\qty{50}{mH}$, si se aplica una corriente que crece regularmente desde cero hasta $I=\qty{10}{A}$ en un tiempo de $t=\qty{0,01}{s}$.

  \item Calcular el coeficiente de autoinducci\'on si se genera una f.e.m de $e=\qty{60}{V}$ con una corriente que va desde $I_i=\qty{0}{A}$ a $I_f=\qty{12}{A}$ en el tiempo de $\Delta t=\qty{0,1}{s}$.

  \item Calcular cu\'anto debe variar la corriente de forma uniforme para que se produzca una f.e.m. de $e=\qty{12}{V}$ en una bobina con coeficiente de autoinducci\'on de $L=\qty{60}{mH}$, en el tiempo de $\Delta t=\qty{0,1}{s}$.

  \item Una bobina posee $N=500$ espiras y produce un flujo magn\'etico de $\Phi=\qty{10}{mWb}$ cuando es atravesada por una corriente de $I=\qty{10}{A}$. Determinar el coeficiente de autoinducci\'on.

  \item Una bobina que posee $N=500$ espiras tiene un coeficiente de autoinducci\'on $L=\qty{0,5}{H}$. Averiguar si produce un flujo magn\'etico de $\Phi=\qty{10}{mWb}$ ¿Cu\'anta corriente circula por la misma?

  \item Averiguar el n\'umero de espiras que tiene una bobina que produce una variaci\'on de flujo magn\'etico de $\Delta \Phi=\qty{0,01}{Wb}$ cuando es atravesada por una corriente de \qty{10}{A} y tiene un coeficiente de autoinducci\'on igual a $L=\qty{0,5}{H}$.

  \item Un conductor de $l=\qty{25}{cm}$ se mueve en forma circular dentro de un campo magn\'etico cortando el flujo magn\'etico seg\'un la funci\'on $\sin(\alpha)$ si el flujo m\'aximo $\Phi=\qty{50}{mWb}$. Si demora $t=\qty{1}{s}$ en dar una vuelta. Realizar un gr\'afico donde el eje de las ordenadas sea el tiempo en segundos y las abscisas el valor de la f.e.m generada.

  \item Un im\'an produce una inducci\'on $\Phi=\qty{100}{mWb}$ en una superficie de $s=\qty{20}{cm^2}$ determinar el valor de la inducci\'on del im\'an.

  \item Un super im\'an de $B=\qty{10}{T}$ tiene secci\'on hexagonal cuyo radio que lo circunscribe es $r=\qty{10}{cm}$. Determinar el flujo magn\'etico en una cara del im\'an $\Phi$.

  \item Un im\'an de $B=\qty{3}{T}$ tiene una longitud de \qty{50}{cm}, se pasa sobre el im\'an un conductor cuya longitud es $l=\qty{20}{cm}$ con una velocidad de $v=\qty{5}{m/s}$. Determinar el valor de la f.e.m. que se produce. ¿A qu\'e velocidad hay que hacerlo pasar para obtener \qty{10}{V}?

  \item Dise\~nar un generador para que genere una f.e.m. \qty{10}{V}, cuando posea un conductor que mueva a la velocidad de \qty{3,14}{m/s}.
\end{enumerate}
\end{document}